% =====================================================================
%  Section: Proposed Research Ideas (Innovative Directions)
% =====================================================================
\section{Proposed Research Ideas}
\label{sec:ideas}

Building on the survey above and the open challenges identified in
Section~\ref{sec:discussion}, we propose four research directions for famous-painting
dynamic video-style generation. All four ideas share a common practical stance: they
build on \textbf{open-source, locally-deployable video diffusion backbones} such as
LTX-Video~2.x and Wan2.1/2.2-I2V, embed \textbf{human artistic guidance} through
hand-drawn drafts (sketches, line drawings, color blocks), provide \textbf{multi-granular
control} over color and stroke distributions, and introduce \textbf{art-creation-domain
innovations} drawn from classical painting theory.

% ---------------------------------------------------------------------
\subsection{Idea 1: Stroke-Anchored Co-Painting Video Generation with Local Video Diffusion}
\label{sec:idea1}

\paragraph{Motivation.} Professional painters do not think in pixels or text prompts;
they think in \emph{strokes}---their direction, pressure, color, and ordering. Current
I2V pipelines accept an image and a prompt, leaving the artist outside the loop.

\paragraph{Core idea.} Build a human-in-the-loop video generation pipeline on top of
Wan2.1-I2V (or LTX-Video) where the artist provides:
(i)~a reference painting (e.g., a Van Gogh canvas) as the style anchor, and
(ii)~a sequence of \emph{hand-drawn stroke drafts}---rough sketches and color blocks
indicating where motion should occur (a swaying cypress, a flowing river). These drafts
act as \emph{spatial--temporal anchors}: a ControlNet-style adapter converts the
stroke drafts into per-frame conditioning, while a stroke-order predictor determines the
temporal evolution of the color blocks.

\paragraph{Control granularity.} The system supports both \textbf{precise} control
(tight line constraints via canny/HED adapters, exact color palette matching via a
histogram-conditional adapter) and \textbf{fuzzy} control (loose color-block regions,
soft attention masks that allow the model to invent texture within a region).

\paragraph{Art-domain innovation.} Inspired by the classical \emph{fat-over-lean}
principle of oil painting, we design a \emph{stroke-temporal curriculum} that schedules
the diffusion denoising steps to first establish broad color composition (lean
underpainting) and then refine high-frequency brush texture (fat impasto layers). This
mirrors how painters physically layer paint, and serves as a domain-informed prior that
improves both temporal coherence and artistic fidelity.

\paragraph{Why it works.} StyleMaster~\cite{ye2025stylemaster} shows that motion adapters
and patch-level style filtering substantially improve stylized video quality; combining
this with stroke-level ControlNet conditioning~\cite{zhang2023controlnet} and a
painting-theoretic curriculum provides a principled way to inject human artistic intent
into a local open-source backbone.

% ---------------------------------------------------------------------
\subsection{Idea 2: Palette-and-Line Dual-Conditioned Ink-Wash Video Synthesis}
\label{sec:idea2}

\paragraph{Motivation.} Chinese ink wash (Shanshui) painting is governed by the
classical \emph{five-level ink theory} (焦、浓、重、淡、清) and deliberate \emph{negative
space} (留白). Existing ink-style methods~\cite{he2023bitgan,liu2026chipdiff,zhang2025inkspirit}
focus on single images and do not expose these art-specific axes as controllable
parameters.

\paragraph{Core idea.} Fine-tune Wan2.1-I2V with a dual-conditional adapter that
accepts:
(i)~a \emph{line-distribution map} (the artist's line draft of mountains, rivers,
pavilions) as a structural prior, and
(ii)~a \emph{color/ink-density distribution map} (hand-painted color blocks specifying
the five-level ink tones and blank-space regions). The two conditions are injected at
different stages of the denoising process: line distribution governs structural
composition in early steps, while ink-density distribution governs tonal rendering in
later steps.

\paragraph{Control granularity.} \textbf{Precise} mode enforces exact ink-density
histogram matching (a five-bin histogram aligned to the five ink levels) for strict
aesthetic adherence; \textbf{fuzzy} mode allows the model to interpolate freely within
each tonal band, producing variations suitable for creative exploration.

\paragraph{Art-domain innovation.} We introduce a \emph{blank-space loss} that explicitly
penalizes texture generation inside user-designated negative-space regions, formalizing
the aesthetic principle of \emph{计白当黑} (treating white as black). Combined with a
\emph{rhythm-of-strokes} prior derived from calligraphic timing, the model is encouraged
to produce temporally coherent ``breathing'' motion (drifting mist, flowing water) that
respects the compositional voids---an effect difficult to achieve with generic video
diffusion priors.

\paragraph{Why it works.} ChipDiff~\cite{liu2026chipdiff} already demonstrates that
loss-gradient guidance on a diffusion model can capture ink aesthetics; extending this
to a video backbone with explicit dual conditioning turns art-theoretic principles into
controllable, locally-runnable tools.

% ---------------------------------------------------------------------
\subsection{Idea 3: Sketch-to-Painting-Video with Hierarchical Fuzzy-to-Precise Control}
\label{sec:idea3}

\paragraph{Motivation.} Artists iterate: a sketch becomes a color study, then a
detailed painting, then an animation. Single-shot I2V generation collapses this
hierarchy into one step, denying the artist intermediate control.

\paragraph{Core idea.} A multi-stage pipeline on LTX-Video where the user supplies a
\emph{progressive sequence of drafts}:
(1)~a rough sketch (line draft),
(2)~a color-block draft (flat regions of color), and
(3)~a detailed reference painting. Each draft activates a different conditional branch:
sketch $\rightarrow$ structure ControlNet, color blocks $\rightarrow$ palette adapter,
reference $\rightarrow$ style LoRA + IP-Adapter. The pipeline generates a short video
clip at each stage, allowing the artist to refine the next draft based on the previous
output---a true \emph{fuzzy-to-precise} co-creation loop.

\paragraph{Control granularity.} A single \emph{control-precision knob} $\lambda \in [0,1]$
interpolates between fuzzy guidance (low $\lambda$, the model freely interprets the
drafts) and precise guidance (high $\lambda$, the output must closely match the
drafts). This is implemented as a learned weighting of conditional features, giving the
artist a continuous dial between creative surprise and faithful execution.

\paragraph{Art-domain innovation.} The staged control mirrors the classical
\emph{sketch $\rightarrow$ underpainting $\rightarrow$ glazing} workflow of Western
atelier practice. We further introduce a \emph{compositional-balance loss} based on the
golden ratio and rule-of-thirds, encouraging drafts and outputs to maintain classical
compositional harmony---a domain prior rarely used in video generation.

\paragraph{Why it works.} Painting Many Pasts~\cite{zhao2020painting} and Stylized Neural
Painting~\cite{zou2021stylized} show that stroke-by-stroke generation is feasible; our
idea lifts this to the video level using a modern open-source diffusion backbone, while
the fuzzy-to-precise dial addresses the controllability gap identified in
StyleMaster~\cite{ye2025stylemaster} and EchoStyle~\cite{li2026echostyle}.

% ---------------------------------------------------------------------
\subsection{Idea 4: Living-Canvas---Collaborative Multi-Region Animation with Artist-Defined Motion Scripts}
\label{sec:idea4}

\paragraph{Motivation.} ``Talking painting'' demonstrations~\cite{zakharov2019fewshot,guo2024liveportrait}
animate a single portrait. But famous paintings are often \emph{canvases with multiple
regions}---a sky, a river, a crowd, a central figure---each of which the artist may want
to animate differently.

\paragraph{Core idea.} A region-aware video generation framework built on
Wan2.1-I2V (multi-region) or LTX-Video. The artist segments the reference painting into
regions and assigns each a \emph{motion script} via hand-drawn arrows, color-flow
blocks, or text. Each region is conditioned independently: a sky region may receive a
soft color-flow block (fuzzy control over cloud drift), while a central portrait region
receives a precise line-drawing of facial motion (precise control for a ``talking
Mona Lisa''). A cross-region consistency module, built on CoDeF-style content
deformation fields~\cite{ouyang2024codef}, ensures seams between regions remain
temporally coherent.

\paragraph{Control granularity.} Per-region $\lambda$ values allow mixed precision: a
precise portrait next to a fuzzy impressionist background. This compositional control is
not possible with monolithic I2V models.

\paragraph{Art-domain innovation.} We introduce the concept of \emph{focal-point
choreography} from classical painting composition: the artist designates a primary
focal region (high precision, strong motion) and supporting regions (low precision,
ambient motion), echoing the artistic principle of \emph{主从虚实} (primary--secondary,
solid--void). The model's attention budget is biased toward focal regions, producing
videos with deliberate artistic hierarchy rather than uniform motion.

\paragraph{Why it works.} RefVFX~\cite{jones2026refvfx} demonstrates reference-guided
temporal effect transfer, and LivePortrait~\cite{guo2024liveportrait} shows efficient
portrait animation; combining region-level conditioning with a local video diffusion
backbone and focal-point choreography yields a system that respects the multi-region
nature of famous paintings while keeping artists in control.

% ---------------------------------------------------------------------
\subsection{Summary of Innovations}

Table~\ref{tab:ideas} summarizes how each idea maps to the four requirements.

\begin{table}[t]
    \centering
    \small
    \begin{tabular}{p{0.20\linewidth} p{0.20\linewidth} p{0.20\linewidth} p{0.20\linewidth} p{0.13\linewidth}}
        \toprule
        Idea & Open Backbone & Human Draft & Control Granularity & Art Innovation \\
        \midrule
        I1: Stroke-Anchored   & Wan2.1 / LTX & Stroke drafts + color blocks & Precise \& fuzzy & Fat-over-lean curriculum \\
        I2: Palette-and-Line  & Wan2.1       & Line + ink-density maps       & Precise \& fuzzy & Five-ink theory, blank-space loss \\
        I3: Sketch-to-Painting-Video & LTX  & Sketch $\rightarrow$ color $\rightarrow$ detail & Fuzzy-to-precise dial & Atelier workflow, golden-ratio loss \\
        I4: Living-Canvas     & Wan2.1 / LTX & Per-region scripts + arrows   & Per-region mixed & Focal-point choreography \\
        \bottomrule
    \end{tabular}
    \caption{Mapping of the four proposed ideas to the design requirements.}
    \label{tab:ideas}
\end{table}

All four ideas are designed to run on a single local GPU using open-weights models,
keep the artist in the loop through hand-drawn drafts at multiple abstraction levels,
offer a continuous spectrum from precise to fuzzy control, and embed principles drawn
from classical painting theory---together forming a coherent research agenda for
famous-painting dynamic video-style generation.
